\documentclass[UTF8,letterpaper,10pt,fontset=none]{ctexart}

\usepackage[letterpaper,left=0.75in,right=0.75in,top=0.72in,bottom=0.72in,columnsep=0.25in]{geometry}
\usepackage{fontspec}
\setmainfont{texgyretermes-regular.otf}[
  BoldFont=texgyretermes-bold.otf,
  ItalicFont=texgyretermes-italic.otf,
  BoldItalicFont=texgyretermes-bolditalic.otf]
\setCJKmainfont[AutoFakeSlant=0.2]{Noto Serif CJK SC}
\setCJKsansfont{Noto Sans CJK SC}
\setCJKmonofont{Noto Sans Mono CJK SC}
\usepackage{caption}
\captionsetup{font=small,labelfont=bf}
\usepackage[unicode,hidelinks]{hyperref}

\ctexset{
  section={format=\normalfont\bfseries\fontsize{11}{13}\selectfont,
    beforeskip=12pt,afterskip=4pt},
  subsection={format=\normalfont\bfseries\fontsize{10}{12}\selectfont,
    beforeskip=9pt,afterskip=3pt},
  subsubsection={format=\normalfont\bfseries\fontsize{10}{12}\selectfont,
    beforeskip=7pt,afterskip=2pt}
}
\setlength{\parindent}{1em}
\setlength{\parskip}{0pt}
\setcounter{secnumdepth}{2}
\pagestyle{empty}
\raggedbottom

\begin{document}
\twocolumn[{
  \vspace*{28pt}
  \begin{center}
    {\fontsize{14}{16}\selectfont\bfseries
      面向 RTL 形式化验证的大语言模型与\\
      证明反馈协同的高效归纳引理选择方法\par}
    \vspace*{105pt}
  \end{center}
  \vspace*{6pt}
}]

\begin{center}\bfseries 摘要\end{center}
\vspace*{125pt}
\noindent\textbf{关键词：}\par
\mbox{}\par

\section{引言}

\subsection{研究背景与动机}

RTL 设计的安全属性验证要求确认：从合法初始状态出发，设计的执行始终满足给定规格。有界模型检查能够发现有限深度内的可达反例，但未发现反例并不能证明任意长度的执行均安全。归纳证明提供了另一条途径。以直接对目标属性进行一步归纳为例，证明既要确认初始状态满足属性，又要确认任意满足属性的当前状态经过一次转移后仍满足属性。后一条件涵盖实际不可达的状态，因此，即使属性在所有可达状态上成立，它也未必能单独构成归纳不变式。辅助引理通过刻画设计状态之间的约束，能够与目标属性一起构成可归纳的证明条件。构造此类引理通常需要理解 RTL 中的状态更新、目标属性与中间状态关系。

大语言模型（LLM）已被用于提出候选引理，再由形式化工具检查其有效性。例如，Peled 等人的 \emph{Large Lemma Miners}~\cite{peled2026large} 将候选生成与形式化筛选结合起来。候选池的形成并不意味着目标属性已经得到证明：候选可能无法解析、不满足初始条件，或虽正确却与目标证明无关；另一些候选则可能只有联合使用才能发挥作用。因此，候选生成之后仍存在独立的选择问题，即应优先检查哪些引理组合。

这一选择问题同时受验证成本约束。不同组合需要分别接受初始性和归纳步骤检查，失败后的再次尝试还会增加形式化查询；若借助 LLM 分析候选或修复组合，模型推理也会消耗墙钟时间。仅减少求解器调用而不计算新增的模型开销，不能说明整体证明过程更高效。本文因此聚焦于固定候选池中的目标导向选择：在保持最终证明由形式化工具确认的前提下，探索如何用候选间的语义联系和证明反馈安排检查顺序，并控制总成本。

\subsection{研究问题与挑战}

本文固定 RTL 设计、目标安全属性、环境假设及候选引理池；选择阶段只从已有候选中取子集，不生成新引理，也不修改待验证的规格。对于被选组合，初始性要求每个合法初始状态都同时满足目标属性和组合中的每条引理；归纳保持则要求任意满足这些条件的状态经过一步合法转移后，后继状态仍满足它们。只有两项均被形式化后端证明，该组合才能用于证明目标属性。这一定义将候选生成能力与选择算法的效果分开，使不同方法可以在相同候选池和形式化后端上比较。

\emph{Large Lemma Miners}~\cite{peled2026large} 的筛选算法先对候选分类排序，再依次检查排序后的前缀。它只访问有限的前缀组合：如果有效组合不是该顺序的前缀，而夹在其中的候选又妨碍归纳，搜索就可能漏过有效组合。经典 Houdini 则从候选全集出发，迭代删除违反归纳条件的候选；在相关查询均完成且语义一致时，它适合求取固定池的归纳子集，但其搜索过程并非专为尽早证明某一指定属性而安排。两者都提供了重要的比较基准，本文不预设其中任何方法在所有任务上更慢或更弱。

目标导向的组合选择面临三项相互关联的挑战。首先，单条引理的检查结果不能完全说明它在组合中的价值：某条引理的归纳步骤失败，可能意味着它需要其他引理支持，而不是应被永久舍弃。其次，形式化失败必须按初始性、目标属性步骤和辅助引理步骤区分；反例状态若无法按照相同的 SVA、时钟和复位语义可靠求值，就不能用于确定性剪枝。最后，形式化检查、LLM 调用与回退搜索共享同一时间预算，需要在可能的证明收益与实际成本之间作出可审计的决策。

据此，本文围绕以下研究问题展开：

\noindent\textbf{RQ1：}在统一预算下，目标导向的 LLM 引导选择能否在维持有竞争力的证明完成率时，减少形式化查询和包含模型开销的证明时间？\par
\noindent\textbf{RQ2：}当候选池扩大并混入无关、冗余或非归纳候选时，这种选择方式的效率和完成率如何变化？\par
\noindent\textbf{RQ3：}组合级初始选择与结构化证明反馈，能否分别带来超出候选排序和仅接收验证成败的重试的增量收益？

\subsection{主要贡献与论文结构}

为研究上述问题，本文提出 LLM-GLS。该方法让 LLM 基于 RTL、目标属性和固定候选池推荐少量候选 ID 组合及排序。形式化后端逐组检查初始性和归纳性，并把可可靠提取的失败义务反馈给选择器。选择器据此修复组合、过滤被已知反例确定否决的组合；成本控制器再根据剩余预算决定继续检查、请求新组合或进入基于完整合格候选池的 Houdini 回退。LLM 只决定检查的优先级，任何成功输出都必须由形式化证明支持。

本文的主要设计与贡献包括：
\begin{enumerate}
  \item 
  \item 
  \item 
\end{enumerate}

后文先介绍相关工作与形式化问题定义，再通过动机示例展示非前缀组合的必要性；随后给出 LLM-GLS 的算法和系统实现，介绍实验设置与结果分析；最后讨论方法的适用边界与研究局限，并在结论中展望后续工作。
% 正式实验完成后，依据实测结果补充主要定量发现，并复核贡献措辞。

\section{相关工作}

\subsection{大语言模型辅助硬件形式化验证}

现有大语言模型辅助硬件形式化验证的研究，首先可以按其生成或处理的对象区分。\emph{AssertLLM} 从设计规格文档提取结构信息并生成待验证的硬件断言\cite{fang2024assertllm}；\emph{FVEval} 则以自然语言到 SVA、设计到 SVA 等任务评估模型生成断言的能力\cite{kang2024fveval}。这类工作主要解决\emph{验证目标如何表达}的问题。另有 \emph{FVDebug} 利用失败轨迹、RTL 和规格定位形式化验证失败的原因，并辅助修改设计\cite{bai2025fvdebug}；其任务是解释或修复失败的实现，而非为既定目标选择证明引理。

与本文更接近的是 \emph{Large Lemma Miners}：它让 LLM 提出 RTL 归纳证明的候选引理，再由形式化后端确认能否构成证明\cite{peled2026large}。该工作同时包含非 Agentic 与 Agentic 候选生成流程；后者会根据候选是否成立等检查结果继续生成引理。因此，\emph{形式化反馈指导 LLM} 已有先例。本文沿用固定的候选生成设置，研究候选池形成之后的组合选择，避免将生成策略的差异计入选择算法的效果。

\subsection{归纳引理生成与不变式推断}

利用验证反馈构造不变式也有较长的研究脉络。\emph{ICE} 用正例、反例和蕴含关系约束不变式学习\cite{garg2014ice}，表明反例不仅能否定候选，还能提供后续搜索所需的结构化约束。在软件验证中，Kamath 等人让 LLM 生成循环不变式，并用符号工具检查\cite{kamath2023loop}；\emph{Lemur} 让 LLM 提议证明子目标，再由自动推理器验证推理步骤\cite{wu2024lemur}。这些工作说明模型提议与形式化判定可以分工，但其主要对象是程序性质或循环不变式的生成，而非给定 RTL 候选池的低成本选择。

在候选组合方面，\emph{Clause2Inv} 让 LLM 生成程序循环不变式的子句，由符号组合器利用历史反例构造完整不变式；违反的证明义务还会引导后续子句生成\cite{cao2025clause2inv}。在反馈修复方面，\emph{LORIS} 检查 LLM 为失败证明义务给出的逐步推理，定位其中的局部逻辑错误，再引导模型重生成循环不变式\cite{li2026loris}。因此，反例驱动组合与结构化反馈并非本文独有的思路。

与上述工作相比，本文的决策对象是生成后冻结的 RTL 引理池。LLM 仅推荐池内引理 ID 的联合组合；形式化后端若返回可可靠求值的初始性或归纳步骤反例，选择器据此过滤或修复\emph{下一组待查组合}。该反馈不产生新引理，也不核查模型的自然语言推理链。模型请求、形式化查询及回退共用墙钟预算。

\subsection{前缀搜索与 Houdini 引理筛选}

在本文给定候选引理池的设置下，\emph{Houdini} 从全部候选出发，依据验证结果迭代剔除不满足初始性或归纳保持条件的候选，直到剩余集合不再变化\cite{flanagan2001houdini}。在候选语义一致且相关查询完成时，它可得到固定池内的最大有效归纳子集。另一种策略是先确定检查顺序：\emph{Large Lemma Miners} 将候选按归纳性与有界检查结果分类排序，再依次验证该顺序产生的前缀\cite{peled2026large}。两种方法采用不同的候选筛选策略。本文在相同候选池、验证条件和预算下，将其作为基线比较证明完成率与检查成本。

后续研究已经把目标信息和搜索成本纳入不变式求解。\emph{H-Houdini} 将硬件不变式学习中的整体归纳检查分解为性质导向、可增量和可复用的层次化检查\cite{dinesh2025hhoudini}；\emph{PyCaliper} 的 RTL 不变式合成框架允许指定候选搜索策略，并在缓存隔离的案例中让 LLM 根据当前不变式和目标性质排列候选顺序\cite{godbole2025pycaliper}。因此，性质导向筛选、增量验证和 LLM 候选排序都已有研究基础。本文关注的是在相同形式化后端和固定候选池上，由 LLM 提议候选 ID 组合，并在可靠的归纳反馈下调整后续待查组合，同时把模型调用与形式化查询纳入同一预算；H-Houdini 的检查分解与并行优化作为相关研究讨论，不作为本文的第三个主外部基线。

\subsection{本文方法定位}

综上，相关工作分别推进了断言或引理生成、反例指导的不变式合成，以及固定候选的排序与筛选。本文选择其中一个更受限的决策点：对同一 RTL 设计、目标属性及已冻结的候选池，预测可能相互支持的非前缀组合；若形式化检查失败，仅在反例语义可靠时据此修复或排除组合，并以统一时间预算决定继续搜索还是回退。形式化后端始终负责初始性与归纳性的最终判定。

本文以相同候选池下的前缀搜索和经典 Houdini 为基线，检验组合选择的收益能否覆盖 LLM 开销；同一次模型输出的排序前缀、以及仅接收成败信息的重试，分别用于检验组合级选择与结构化反馈的增量作用。

\section{预备知识与问题定义}

\subsection{RTL 安全属性与转移系统}

给定 RTL 设计 $D$、目标安全属性及环境假设，形式化后端依据统一的时钟、复位和 SVA 语义构造状态转移系统 $T=(V,\mathrm{Init},\mathrm{Tr})$。其中，$V$ 表示当前状态变量，$V'$ 表示其后继状态变量；$\mathrm{Init}(V)$ 刻画合法初始状态，$\mathrm{Tr}(V,V')$ 刻画合法的一步转移。本文用状态谓词 $p(V)$ 表示展开后的目标安全条件，并把从初始状态经有限次合法转移可到达的状态集合记为 $\mathrm{Reach}(T)$。证明目标是 $T\models AGp$，即对每个 $s\in\mathrm{Reach}(T)$ 均有 $p(s)$ 成立。对含时序运算的属性，上述状态记号指后端按相同语义展开后的模型，而不是直接对原始 SVA 文本作布尔求值~\cite{peled2026large}。

有界模型检查（BMC）从合法初始状态搜索长度不超过给定深度 $d$ 的执行轨迹，可用于发现违反 $p$ 的可达反例。未找到反例只能说明该深度内未发现目标违例，不能据此报告 \texttt{PROVED}；证明成功仍须依靠下文所述的初始性与归纳保持检查。

\subsection{归纳增强与形式化检查}

状态谓词 $I(V)$ 是归纳不变式，当且仅当它满足初始性和归纳保持：
\[
\begin{array}{l}
\mathrm{Init}(V)\Rightarrow I(V),\\
I(V)\land\mathrm{Tr}(V,V')\Rightarrow I(V').
\end{array}
\]
若还满足 $I(V)\Rightarrow p(V)$，它便能证明目标安全属性。目标 $p$ 本身可能不满足归纳保持；此时可用辅助引理组成归纳增强。对任意候选集合 $S$，定义
\[
I_S(V)=p(V)\land\bigwedge_{\ell\in S}\ell(V),
\]
其中空集合上的合取取真，故 $I_{\emptyset}=p$。所选引理无需各自构成归纳不变式；只有整个 $I_S$ 满足上述两项条件，$S$ 才是目标 $p$ 的有效归纳增强。

实际检查可表述为两项可满足性查询：
\[
\begin{array}{rl}
Q_{\mathrm{init}}(S)={}&\mathrm{Init}(V)\land\neg I_S(V),\\
Q_{\mathrm{step}}(S)={}&I_S(V)\land\mathrm{Tr}(V,V')\\
&{}\land\neg I_S(V').
\end{array}
\]
两项均为 UNSAT 才能返回证明；$Q_{\mathrm{step}}$ 的 SAT 见证可能始于不可达状态，不能直接视为目标属性的可达反例。$Q_{\mathrm{init}}$ 为 SAT 时也须区分是 $p$ 还是某条候选在初始状态失败。TIMEOUT、UNKNOWN 或工具错误均不构成证明。对无法按后端语义可靠判定单步反例中候选真值的时序 SVA，反馈求值记为 unknown，不据此作确定性剪枝。

\subsection{固定候选池上的引理选择问题}

候选生成结束后，固定引理池 $C=\{\ell_1,\ldots,\ell_n\}$，并为每条候选分配不变的 ID；候选内容在选择阶段保持不变。给定同一设计 $D$、目标 $p$、环境假设、后端语义及候选池 $C$，选择器安排待检验子集 $S\subseteq C$ 的顺序；它不生成新引理，也不修改设计或目标属性。候选仅表示待验证的性质，并不预设其满足初始性、归纳保持或与目标相关。下文公式中的 $S$ 表示引理集合，实际输出则由其 ID 列表唯一表示。

固定顺序的前缀搜索只访问以下组合：
\[
P_k=\{\ell_1,\ldots,\ell_k\},\qquad 1\le k\le n.
\]
一般的选择问题还允许非前缀组合。形式化目标是找到某个 $S\subseteq C$，使 $Q_{\mathrm{init}}(S)$ 与 $Q_{\mathrm{step}}(S)$ 均为 UNSAT，而不是预先要求最少的引理数。失败的初始状态或步骤状态对可为后续组合选择提供信息，但任何据此提出的新组合都必须重新接受完整的形式化检查；若后端不能可靠解释见证中的候选取值，则不能把该见证当作安全剪枝依据。

\subsection{评价目标与结果状态}

本文在固定候选池就绪后，为每个任务设置覆盖选择与证明阶段的统一墙钟预算 $B$。LLM 调用、全部形式化查询、必要的反例求值和回退搜索均消耗该预算；候选生成是各方法共用的前置阶段，不计入选择方法的 Proof Time。评价首先比较预算内完成证明的目标属性数量（Solved Properties），再比较各方法共同成功任务从候选池就绪到证明成立的墙钟时间（Proof Time）与形式化查询次数（Formal Queries）。Proof Time 包含模型推理、形式化检查及其他选择开销；Formal Queries 按实际发起的后端请求计数，包含为定位失败或求值反例而追加的查询。Solver Time、LLM Calls／Tokens 和所选引理数作为辅助指标。未证明的任务保留其终止状态与已消耗时间，并计入完成率的分母；其已耗时间不作为成功证明时间纳入上述成本均值。

结果状态须反映证明证据的强度：\texttt{PROVED} 表示某个组合的两项归纳查询均已被证明 UNSAT。若发现从初始状态可达的 $\neg p$ 反例，则目标属性被违反；归纳步骤的 SAT 见证只否决当前组合，不能单独说明目标属性被违反。

\texttt{NO\_PROOF\_IN\_POOL} 仅在固定池搜索完整且可靠、又无有效组合时返回；它不表示目标属性不安全。预算耗尽、检查未决及工具故障分别记录为 \texttt{TIMEOUT}、\texttt{UNKNOWN} 与 \texttt{ERROR}，不得与证明成功或可达反例混同。

\section{动机示例}

考虑 Large Lemma Miners 开源基准中的 \texttt{swap2}~\cite{peled2026artifact}。设计中三个寄存器的初始条件为 $\mathrm{Init}:(a,b,c)=(1,5,7)$，转移关系 $T$ 由 $a'=b$、$b'=c$、$c'=a$ 给出。令 $K=\{1,5,7\}$，待证明的目标属性为 $p: a\in K$。所有可达状态中，$a,b,c$ 始终是 $K$ 中三个值的排列，因此 $p$ 成立；但单独的 $p$ 不是一步归纳的：例如不可达的源状态 $(1,0,0)$ 满足 $p$，经过一步转移后 $a'=0$，违反 $p$。

为说明候选选择，设有三条在可达状态上均成立的候选引理，并按示意顺序排列：$\ell_1:a\ne b$、$\ell_2:b\in K$、$\ell_3:c\in K$。考虑组合 $S^*=\{\ell_2,\ell_3\}$：它不是该顺序的前缀。三个前缀 $\{\ell_1\}$、$\{\ell_1,\ell_2\}$、$\{\ell_1,\ell_2,\ell_3\}$ 分别可由源状态 $(1,0,0)$、$(1,5,0)$、$(1,5,5)$ 构造一步反例：相应后继依次违反 $p$、$\ell_2$、$\ell_1$。因此，只检查这些前缀无法检验到 $S^*$。下面用一步归纳检查验证这个非前缀组合，并展示反例如何引导组合修复。

下列检查中，$I_0$ 和 $I^*$ 分别表示目标属性与相应候选组合的合取；对任意状态谓词 $I$，$I'$ 表示将其当前状态变量替换为后继状态变量得到的谓词。

\noindent\begin{minipage}{\columnwidth}
\hrule\smallskip
\noindent\textbf{示例检查：从 $S_0$ 到 $S^*$}\par
\smallskip\hrule\smallskip
{\small
\begin{enumerate}
  \setlength{\topsep}{2pt}
  \setlength{\itemsep}{1pt}
  \setlength{\parsep}{0pt}
  \item 取 $S_0=\{\ell_2\}$，令 $I_0=p\land\ell_2$。$\mathrm{Init}\land\neg I_0$ 不可满足，故初始性成立。
  \item 步骤查询 $I_0\land T\land\neg I_0'$ 可满足：$(1,5,0)\to(5,0,1)$。后继中 $b'=0\notin K$，故归纳保持失败。
  \item 由 $b'=c$ 可知，$\ell_2'$ 失败源于源状态的 $c=0$。加入 $\ell_3:c\in K$ 排除此状态，令 $I^*=p\land\ell_2\land\ell_3$。
  \item 将前述两项查询中的 $I_0$ 替换为 $I^*$，结果均不可满足；故 $I^*$ 是归纳不变式，且 $I^*\Rightarrow p$。
\end{enumerate}
}
\hrule
\end{minipage}\par\smallskip

这说明候选组合的归纳步骤失败未必应导致全局删除其中的引理；它也可能提示需要另一条候选提供联合支持。

\section{LLM-GLS：证明反馈驱动的归纳引理组合选择}
\mbox{}\par
\subsection{总体架构}
\mbox{}\par
\subsection{候选预处理与初始组合提议}
\mbox{}\par
\subsection{形式化反馈与三值反例求值}
\mbox{}\par
\subsection{安全剪枝与局部修复}
\mbox{}\par
\subsection{成本感知排序与预算决策}
\mbox{}\par
\subsection{Houdini 回退与正确性边界}
\mbox{}\par
\subsection{完整算法流程}
\mbox{}\par

\section{实验设置}
\subsection{实验任务与候选池}
\subsubsection{固定候选池数据集}
本文采用 \emph{Large Lemma Miners} 公开的 RTL/SVA 基准任务及其可复现工件\cite{peled2026large,peled2026artifact}。每个任务由一个 SystemVerilog 设计和一条待证安全属性构成。沿用该基准的任务划分，主实验组包含 78 个任务，困难实验组包含 31 个任务，共 109 个任务。原工作将至少一种对比模型检查器未能在 10 分钟内证明目标属性的任务归为需要非平凡引理的困难案例\cite{peled2026large}；其余任务构成主实验组。该划分在本文的方法比较前确定，不依据 PS、CH 或 LLM-GLS 的运行结果调整。

为隔离引理选择策略的影响，本文在正式比较前为每个任务生成并冻结一个候选引理池。两个实验组使用相同的非 Agentic 生成配置：模型为 \texttt{gpt-6.1-sol}，推理强度为 \texttt{xhigh}，提供一个 few-shot 示例并采样五次。生成的候选按文本完全一致去重，保留首次出现的顺序并分配候选 ID；任务、属性和逐任务候选池均记录哈希以供核对。PS、CH 和 LLM-GLS 对同一任务读取完全相同的冻结候选池，选择阶段不再生成新候选。

\subsubsection{扩展候选池数据集}
\mbox{}\par
\subsection{对比基线}
\mbox{}\par
\subsection{评价指标与统计方法}
三种方法共用预先生成并冻结的候选池。候选生成耗时不参与选择方法的证明时间比较；证明时间包含选择阶段的 LLM 推理、形式化查询及其他处理开销。结果分析分别报告模型与形式化后端的实际耗时，以检验新增的模型开销是否被节省的验证开销抵消。
\subsection{公平性设置与参数配置}
三种方法采用统一的工具错误终止策略。任一形式化查询因编译错误、后端不支持、进程崩溃或启动失败返回 \texttt{ERROR} 时，当前任务立即结束，并保存错误阶段、候选组合、退出码及后端日志；批次继续处理下一任务。该规则覆盖预检查、候选分类、组合验证、见证提取和 Houdini 搜索，工具错误不触发额外的模型修复。有效反例、\texttt{TIMEOUT} 与 \texttt{UNKNOWN} 分别沿用反例修复及预算管理规则，不与工具错误混同。

\section{实验结果与分析}
\mbox{}\par
\subsection{实验一：证明有效性与效率}
\subsubsection{证明完成率}
为评价不同引理选择方法的证明有效性，本文分别在包含 78 个任务的主实验组和包含 31 个任务的困难实验组上，比较 Prefix Search（PS）、Classical Houdini（CH）和 LLM-GLS 的证明完成情况。各方法使用相同的任务、固定候选引理池、形式化后端与时间预算。只有目标属性及所选引理的初始性和归纳性均得到形式化证明、最终返回 \texttt{PROVED} 的任务计入成功；\texttt{TIMEOUT}、\texttt{UNKNOWN} 和 \texttt{ERROR} 分别统计，并保留在任务总数中。方法 $M$ 在实验组 $G$ 上的证明完成率定义为
\[
  \mathrm{CR}(M,G)=
  \frac{N_{\mathrm{PROVED}}(M,G)}{N_G}\times 100\%,
\]
其中 $N_G$ 为该组任务数。实验一对各方法独立重复运行三轮，证明完成率与后续证明成本分析均使用这同一批运行记录。完成率按轮对全部任务计算，再报告三轮的均值与标准差，避免将重复运行视为新增任务。

\begin{table}[t]
  \centering
  \caption{主实验组与困难实验组的证明完成情况}
  \label{tab:proof_completion_rate}
  \small
  \setlength{\tabcolsep}{4pt}
  \begin{tabular}{lrrrrr}
    \hline
    方法 & 成功 & 超时 & 未决 & 错误 & 完成率（\%） \\
    \hline
    \multicolumn{6}{l}{主实验组（78 个任务）} \\
    PS      & 62 & 0 & 14 & 2 & 79.49 $\pm$ 0.00 \\
    CH      & 66 & 0 & 10 & 2 & 84.62 $\pm$ 0.00 \\
    LLM-GLS & 71 & 0 & 5  & 2 & 91.03 $\pm$ 0.00 \\
    \hline
    \multicolumn{6}{l}{困难实验组（31 个任务）} \\
    PS      & 20 & 2 & 3 & 6 & 64.52 $\pm$ 0.00 \\
    CH      & 20 & 2 & 3 & 6 & 64.52 $\pm$ 0.00 \\
    LLM-GLS & 23 & 1 & 1 & 6 & 74.19 $\pm$ 0.00 \\
    \hline
  \end{tabular}
  \par\smallskip
  \begin{minipage}{\columnwidth}
    \footnotesize
    注：成功、超时、未决、错误分别对应 \texttt{PROVED}、\texttt{TIMEOUT}、\texttt{UNKNOWN}、\texttt{ERROR}。状态数量为三轮运行的均值，完成率以三轮均值 $\pm$ 标准差表示。
  \end{minipage}
\end{table}

表~\ref{tab:proof_completion_rate} 汇总了两组实验结果。三轮运行的完成率标准差均为 $0.00$ 个百分点。在主实验组中，PS、CH 和 LLM-GLS 分别完成了 62、66 和 71 个任务，证明完成率分别为 79.49\%、84.62\% 和 91.03\%。LLM-GLS 相较 PS 和 CH 分别提高了 11.54 和 6.41 个百分点。这表明在主实验组和给定预算下，LLM-GLS 的完整流程比两种基线完成了更多证明任务。

在困难实验组中，PS 和 CH 均完成了 20 个任务，完成率均为 64.52\%；LLM-GLS 完成了 23 个任务，完成率为 74.19\%，比两种基线均多完成 3 个任务，按未四舍五入的完成率计算提高 9.68 个百分点。与主实验组相比，三种方法在困难实验组的完成率均有所下降，但 LLM-GLS 在该组仍保持最高完成率。这一结果表明，LLM-GLS 的完成率优势在预先划定的困难实验组中也得到观察。

按三轮运行的均值统计未完成任务的终止状态，主实验组中，PS、CH 和 LLM-GLS 分别有 14、10 和 5 个任务返回 \texttt{UNKNOWN}，各有 2 个任务返回 \texttt{ERROR}，均无 \texttt{TIMEOUT}。困难实验组中，PS 和 CH 各有 2 个 \texttt{TIMEOUT}、3 个 \texttt{UNKNOWN} 和 6 个 \texttt{ERROR}；LLM-GLS 有 1 个 \texttt{TIMEOUT}、1 个 \texttt{UNKNOWN} 和 6 个 \texttt{ERROR}。由于同组各方法的 \texttt{ERROR} 数量相同，LLM-GLS 在主实验组增加的证明成功数对应更少的 \texttt{UNKNOWN}，在困难实验组则对应更少的 \texttt{TIMEOUT} 和 \texttt{UNKNOWN}。这些状态区分了预算不足、证明未决和工具执行异常，不能统一解释为目标属性不成立或候选池中不存在有效引理组合。因此，上述完成率反映的是完整流程在给定预算下的证明覆盖范围。下一小节在三种方法共同成功的任务上比较证明时间与形式化查询量，以避免不同成功任务构成影响成本比较。

\subsubsection{证明时间与形式化查询量}
在三种方法均完成证明的 78 个任务中，主实验组有 58 个，困难实验组有 20 个。证明时间包含选择、形式化检查及模型推理等实际开销，候选生成时间作为共用前置成本单独处理，不作为对比的指标。表~\ref{tab:proof_time_queries} 汇总了各方法在这些任务对应的 234 条记录上的均值。LLM-GLS 的平均证明时间为 9.70 s，相较 PS 的 37.88 s 和 CH 的 11.23 s，分别减少 74.39\% 和 13.61\%。

\begin{table}[t]
  \centering
  \caption{共同成功任务上的证明成本比较}
  \label{tab:proof_time_queries}
  \small
  \setlength{\tabcolsep}{4pt}
  \begin{tabular}{lrrr}
    \hline
    方法 & 证明时间（s） & 查询次数 & LLM 调用次数\\
    \hline
    PS      & 37.88 & 20.46 & 0\\
    CH      & 11.23 & 7.72 & 0\\
    LLM-GLS & 9.70  & 5.68 & 1.13\\
    \hline
  \end{tabular}
  \par\smallskip
  \begin{minipage}{\columnwidth}
    \footnotesize
    注：各项为三轮运行记录的均值；每种方法包含相同的 78 个任务、共 234 条记录。
  \end{minipage}
\end{table}

在同一任务集合上，LLM-GLS 的平均形式化查询次数为 5.68 次，相较 PS 的 20.46 次和 CH 的 7.72 次，分别减少 72.24\% 和 26.41\%。在这些共同成功任务上，LLM-GLS 平均调用 1.13 次 LLM。

由于证明时间已经包含模型推理，表~\ref{tab:proof_time_queries} 中较短的平均证明时间表明，新增模型调用没有抵消这一任务集合上的总体时间收益。

综合表1的证明完成率，以及表2在共同成功任务上的平均证明时间和形式化查询次数，实验一表明：在固定候选池与统一预算下，LLM-GLS 在主实验组和困难实验组均完成更多证明任务；

在三种方法共同成功的任务上，其平均证明时间更短、形式化查询更少。非前缀组合选择与结构化证明反馈各自的贡献，留待后续关键设计对照实验分析。

\subsection{实验二：候选池规模与干扰鲁棒性}
\mbox{}\par
\subsubsection{候选池规模的影响}
\mbox{}\par
\subsubsection{干扰候选类型的影响}
\mbox{}\par
\subsection{实验三：关键设计对照}
\mbox{}\par
\subsubsection{推荐组合与排序前缀对照}
\mbox{}\par
\subsubsection{结构化证明反馈对照}
\mbox{}\par
\subsubsection{修复轮次与调用成本}
\mbox{}\par
\subsection{典型任务与未决案例分析}
\mbox{}\par

\section{讨论与研究局限}
\mbox{}\par
\subsection{方法适用范围}
\mbox{}\par
\subsection{形式化证据与语义边界}
\mbox{}\par
\subsection{数据集与候选池偏差}
\mbox{}\par
\subsection{模型随机性与成本测量}
\mbox{}\par

\section{结论与未来工作}
\mbox{}\par


\section*{致谢}

\bibliographystyle{unsrt}
\begingroup
\raggedright
\bibliography{references}
\endgroup

\appendix
\section{算法与输出契约}
\mbox{}\par
\subsection{完整算法伪代码}
\mbox{}\par
\subsection{候选 ID 与组合输出格式}
\mbox{}\par
\subsection{三值求值与安全剪枝条件}
\mbox{}\par

\section{实验复现配置}
\mbox{}\par
\subsection{任务清单与候选池哈希}
\mbox{}\par
\subsection{提示词与模型配置}
\mbox{}\par
\subsection{形式化后端与预算参数}
\mbox{}\par
\subsection{成本参数的校准与冻结}
\mbox{}\par

\section{补充实验与任务级结果}
\mbox{}\par
\subsection{完整任务级结果}
\mbox{}\par
\subsection{参数敏感性分析}
\mbox{}\par
\subsection{失败与未决案例}
\mbox{}\par


\end{document}
